% BLAÐSÍÐU- OG LETURSTÆRÐ, SPÁSSÍUR, MÁLSGREINAR
\documentclass[12pt]{article}    % Hægt að nota 10 pt, 11 pt, report, book
\usepackage[a4paper]{geometry}   % Eða t.d. a3paper, letterpaper
  \geometry{margin=2.2cm}        % Eða left, right, top, bottom, 
  % vmargin, hmargin
\usepackage[skip=0.8em]{parskip} % Bil milli málsgreina, og ekki inndregnar

% FYRIR FONTA OG ÍSLENSKU
\usepackage{fourier}          % Fallegri fontar en sá sjálfgefni, sér í
                              % lagi í formúlum (aðrir möguleikar eru 
                              % \usepackage[nott]{kpfonts}) og fyrir
                              % sans-serif letur \usepackage{cmbright}
\usepackage[utf8]{inputenc}   % Til íslenska o.fl. í inntaksskrá lesist rétt
\usepackage[T1]{fontenc}      % Til að íslenska o.fl. prentist rétt í pdf-skrá
\usepackage[icelandic]{babel} % Gefur t.d. "Mynd xx" undir myndir og skiptir 
                              % orðum rétt milli lína. Telja má upp fleiri
                              % tungumál, með kommum á milli og þá er skipt
                              % á milli í skjalinu með \setlanguage{mál}
\usepackage{microtype}        % betri bil milli stafa, m.a. jafnari hægri brún

% LITIR
\usepackage[svgnames]{xcolor} % Gefur sömu litanöfn og HTML, css og Matplotlib.
                              % Með [x11names] má fá enn fleiri liti, sbr. R 
                              % (ggplot2). Litina má skoða í kafla 4.3 í xcolor 
                              % handbókinni (gúglið xcolor pdf)

% FYRIRSAGNIR
\usepackage{titlesec}                      % Stjórn kaflafyrirsagna
\titlespacing{\section}{0ex}{*3}{*1}       % Bil framan kaflanúmers 0pt, ofan 
\titlespacing{\subsection}{0ex}{*3}{*1}    % fyrirsagna 3ex og neðan þeirra 1ex
\titlespacing{\subsubsection}{0ex}{*3}{*1} % með smá teygju (1ex ≈ 0.4em)
\titlelabel{\thetitle\hspace{0.6em}}       % Bil á eftir kaflanúmeri
% ATH. Sjálfgefið er {0ex}{*3.5}{*2.3} fyrir \section, {0ex}{*3.25}{*1.5}
% fyrir hinar tvær og 1em fyrir bil á eftir. Með \usepackage[compact]{titlesec} 
% verður sjálfgefið {*1}{*1} fyrir allar þrjár. Svo má stilla font og lit 
% fyrirsagnar, t.d. með \titleformat*{\section}{\Large\bfseries\color{DarkBlue}}

% MYNDIR OG FORRIT
\usepackage{graphicx}         % Fyrir myndir
\usepackage{minted}           % Fyrir forrit. Nokkrir (ekki of litríkir) stílar
  \usemintedstyle{friendly}   % eru vs, friendly, tango, trac, borland, bw

% STÆRÐFRÆÐIFORMÚLUR
\usepackage[leqno]{mathtools} % Bætir mörgu gagnlegu við stærðfræðiformúlur
\usepackage{xfrac}            % Svo má nota sfrac til að fá 1/2, 2/3 o.s.frv.
\usepackage{nicefrac}         % Aðeins öðruvísi brot

% BETRI TÖFLUR OG LISTAR
\usepackage{booktabs}         % Fyrir betri töflur
\usepackage{caption}          % M.a. aukabil á undan töflufyrirsögnum
  \captionsetup{font=small,   % Smærra letur á mynda- og töflutexta,
                labelfont=bf, % "Mynd 1" og "Tafla 2" feitletrað 
                margin=2cm}   % Viðbótarspássíur á mynda- og töflutexta
\usepackage{siunitx}          % Virkjar m.a. S-jöfnun í töflum
\sisetup{output-decimal-marker={,}} % Tugabrot með kommu, t.d. 5,37 í stað 5.37
\usepackage[shortlabels]{enumitem} % Fyrir betri lista, t.d. a) b)...
  \setlist{nosep,             % Einfalt línubil í listum  
          topsep=-.5\parskip} % Það er smálús í samtvinnun parskip og enumitem

% TENGLAR          x<x  
\usepackage{hyperref}         % Fyrir tengla
  \hypersetup{colorlinks,     % Ef sleppt þá allir tenglar svartir
            linkcolor=Indigo, % tilvísun innan skjals
            citecolor=Indigo, % – í heimildaskráhttps://www.overleaf.com/project/607851adb12f15271592fea3
            urlcolor=Indigo}  % – út á netið

\author{Setjið inn höfund hér}
\date{Setjið inn mánuð og ár hér}

% ================= HÉR ENDAR SNIÐMÁTIÐ =======================

% Gefa dæmi um skilgreiningu nýrrar skipunar:
\newcommand{\rautt}[1]{\textcolor{FireBrick}{#1}}  % [1] er fjöldi stika
 
% Sýna leið til að breyta leturgerð og -stærð og línubili í fyrirsögn skjals:
\title{\textbf{\Huge Latex sýnishorn}\vspace{-1em}} 
\author{Kristján Jónasson\\Stærðfræði og reiknifræði}
\date{Vorið 2021}

\begin{document}

\maketitle

\section{Um alheiminn}  % Með \documentclass{report} notar maður líka \chapter
Til er kenning sem segir að ef einhver muni uppgötva námkvæmlega til hvers heimurinn er og af hverju hann er hér þá muni hann hverfa samstundis og í staðinn komi eitthvað enn skrýtnara. Til er önnur kenning sem segir að þetta hafi þegar gerst.\cite{adams1995hitchhiker}

Hér kemur svo ný málsgrein.
%
\begin{figure}[h!]  % h=helst hér (here) og h!=endilega hér. t=top, b=bottom
    \centering
    \includegraphics[width=0.8\textwidth]{eldgos.jpg}
    \caption{Eldgosið}
    \label{mynd:eldgos}
\end{figure}

\section{Latex-sýnishorn}
\subsection{Ýmis atriði}
\textit{Skáletur}, \textbf{feitletur} \texttt{ritvélaletur}, \textcolor{MediumBlue}{blátt letur} og \rautt{rautt letur}. Hér er ,,texti innan gæsalappa``. Sérstök Latex tákn má setja inn (með ritvélaletri) með \texttt{verb} svona: \verb:&%$#_{}~^\: og fyrir sjö fyrstu má nota \verb:\:, t.d. \$300, 50\%, Brauð \& Co. Hér er dæmi þar sem skipta þarf orði milli lína því það er svo langt: vaðlaheiðarvegavinnuverkfærageymsluskúraútidyralykill. Hér eru tilvísanir í mynd, kafla og jöfnu: sjá mynd \ref{mynd:eldgos}, kafla \ref{kafli:nidurstada} og jöfnu (\ref{eq:rpol}). Og hér er tengill: \href{https://en.wikipedia.org/wiki/Universe}{Alheimurinn á Wikpediu}, og url: \url{http://hi.is}.

Langt strik: Á árunum 2013--2017 gerðist ýmislegt.
Formúlur: $x_3 = \sqrt{x_1^2 + x_2^2}$ og á sér línu:
%
$$  % Takið eftir % á línunni að ofan. Alveg auð lína mundi gefa greinaskil
\text{Rúmmál kúlu} = \frac{4\pi r^3}{3}
$$  % oft er notað \[ og \] utan um svona jöfnur í stað $$ og $$

\subsection{Listar}
Kúlulisti (bullet list):
%https://www.overleaf.com/project/607851adb12f15271592fea3
\begin{itemize}
    \item Fyrsta atriði
    \item Annað atriði
\end{itemize}
%
tölusettur listi:
\begin{enumerate}
    \item Fyrsta atriði
    \item Annað atriði
    \begin{enumerate}[1.]
        \item undiratriði 1
        \item undiratriði 2
    \end{enumerate}
\end{enumerate}
%
og það má líka merkja með bókstöfum:
%
\begin{enumerate}[a)]
    \item Fyrsta atriði
    \item Annað atriði
\end{enumerate}

\subsection{Forrit}
Hér er dæmi um listun forrits með pakkanum \textit{minted}, en eldri pakki sem litar ekki, dugar ekki á utf-8 stafi og notar öðruvísi (og e.t.v. verri) font er \textit{listings}.
% Ath. Minted pakkinn dugar í Overleaf en á eigin tölvu þarf að maður að
% vera með bæði Python og Pygments (sem fylgir oft með Python), og svo
% þarf -shell-escape að vera með þegar latex eða pdflatex skipun er gefin.
% Svo má fá forrit án "syntax highlighting" með verbatim.
\begin{quote}           % dregur inn línur
\begin{minted}{python}
def gera_tvisvar(f):
    '''framkvæmir f tvisvar'''
    f()
    f()
def bulla():
    '''skrifar "bull"'''
    print("bull")
b = bulla         # b er fallbreyta
gera_tvisvar(b)   # prentar bull tvisvar
\end{minted}
\end{quote}

\subsection{Ljóð}
\begin{verse}
\textrm{RAUN VÍSINDA} \\[.3em]  % 1 em er u.þ.b. hæð textalínu
Raun vísinda:\\
Stofnun háskólans.\\
\hspace{4em}{\footnotesize Dagur Sigurðarson}
\end{verse}

\subsection{Tafla}
Hér er dæmi um töflu með fjórum dálkum, vinstri jöfnuðum, miðjuðum, kommujöfnuðum og miðjuðum sem blokk og loks hægri jöfnuðum.
%
\begin{table}[ht!]  % hér! sbr. figure að ofan
    \centering
    \caption{Gengi 8. apríl 2021. Upplýsingarnar í þessari töflu eru fengnar af vef Landsbankans og vef Seðlabanka Íslands.}
    \sisetup{group-separator={.}} % bil er sjálfgefið
    \begin{tabular}{lcS[table-format=3.3]r} % 3.3 = max ###.###
        \toprule  % Ath slaufusviga utan um titil á S-dálki
        Gjaldmiðill & Skst & {Meðalgengi} & Sölugengi \\
        \midrule
        Króna         & ISK  & 1    & 1000 \\
        Evra          & EUR  & 151,9  & 6,57 \\
        Sterlingspund & GBP  & 175,29 & 5,69 \\
        Dollari       & USD  & 127,79 & 7,82 \\
        Dönsk króna   & DKK  & 20,421 & 32\,248,86 \\
        \bottomrule
    \end{tabular}
\end{table}
%
\section{Fleiri stærðfræðiatriði}
\subsection{Leiðir til að setja jöfnur á sér línur} 
Tölusettar :
%
\begin{equation} \label{eq:pytha}
    a^2 + b^2 = c^2
\end{equation}
%
Fleiri jöfnur saman, miðjaðar (jafna (\ref{eq:rpol}) er afleiðing af jöfnu (\ref{eq:pytha})):
%
\begin{gather}
    \label{eq:rpol}
    r = \sqrt{x^2 + y^2} \\               % \\ til að fara í næstu línu
    \theta = \operatorname{arctan2}(y,x)  % latex þekkir ekki atan2
\end{gather}
%
Fleiri jöfnur saman, standast á við =-merki:
\begin{align*}  % stjarnan til að sleppa tölusetningu 
    r &= \sqrt{x^2 + y^2} \\
    \theta &= \operatorname{arctan2}(y,x)
\end{align*}  % líka má nota * með gather og equation
%
\subsection{Summur og fleira}
Summa:
$$
\sum_{k=1}^m k=\frac{m(m+1)}{2}
$$
Heildi og ýmis (misflókin og misflott) brot:
$$
\int_0^1 x^2 dx = \nicefrac{1}{3} = \text{\large{\sfrac{1}{3}}} 
= \tfrac{1}{3} = 1\big/3 = \frac{1}{3}
$$
Markgildi:
$$
\lim_{x \to 0}\frac{\sin x}{x} = 1  % Án \ túlkast sin sem margfeldið s·i·n
$$
Mengjatákn:
\begin{gather*}
x \in (A \cup B) \cap C 
\text{ og } y \in \big\{1,2,3,4\big\} \\
M \subseteq\mathbb{R}^+ \text{ og } M \neq \emptyset
\end{gather*}
Gríska stafrófið:
\begin{gather*}
\alpha, \beta, \gamma, \delta, \epsilon, \varepsilon, \zeta, \eta, \theta, \iota, \kappa, \lambda, \mu, \nu, \xi, o, \pi, \rho, \sigma, \tau, \upsilon, \theta, \phi, \varphi, \chi, \psi, \omega \\
\Gamma, \Delta, \Theta, \Lambda, \Xi, \Pi, \Sigma, \Phi, \Psi, \Omega.
\end{gather*}

\subsection{Fylki og vigrar}
$$
Ax = \begin{pmatrix}  % p fyrir parentheses
1 & 2\\
3 & 4
\end{pmatrix}
\begin{pmatrix} 5\\6 \end{pmatrix}
= \begin{bmatrix}  % b fyrir brackets
1 & 2\\
3 & 4
\end{bmatrix}
\begin{bmatrix} 5\\6 \end{bmatrix}
$$

\section{Niðurstaða} % Sýnir hver m.a. hvernig stinga má inn
% inn enskum texta, og enskar gæsalappir.
\label{kafli:nidurstada}
\begin{otherlanguage}{english}
``I always thought something was fundamentally wrong with the universe'' \cite{adams1995hitchhiker}. Internationalization is a long word that needs to be hyphenated.
\end{otherlanguage}

%% HEIMILDASKRÁ
\bibliographystyle{plain}
\bibliography{heimildir.bib}
\end{document}